\scanpage{250}
Voici un résultat \deleted{\unsure{plus}} profond, fondamental
\deleted{\unsure{dans la th}} en géométrie analytique :

\uline{Lemme} 1.3. (Grauert). Soit $f:X\unsure{-|\to}Y$ un
morphisme propre d'espaces analytiques sur $\mathbb C$,
$F$ un faisceau coh. sur $X$. Alors

(i) les $R^if_{\unsure{*}}(\deleted{\unsure{E}}F)$ sont cohérents

(ii) Soit $y\in Y$, \deleted{\unsure{Alors le système}}\inserted{$Y_n=\unsure{\operatorname{Specan}|\operatorname{Spec\,an}}(\mathcal O_y/m_y^{n+1})$},
$X_n=X\times_Y Y_{\unsure{n|m}}$, $F_n=F\otimes_Y Y_n$, alors le système
projectif $(H^i(X_n,F_n))_n$ satisfait M.L.
et l'hom canonique
\[
R^if_*(F)_y{}^{\wedge}\longrightarrow\underset{n}{\underleftarrow{\lim}}H^i(X_n,F_n)
\]
est un isom .

\uline{Corollaire} 1.4. Supposons que
$y$ soit un point isolé de $\operatorname{supp} R^if_*(\unsure{E|F})$,
\deleted{Alors} i.e que $R^if_*(F)_y$ soit artinien.
Alors on a un isom
\[
R^if_*(F)_y\longrightarrow\underset{n}{\underleftarrow{\lim}}H^i(X_n,F_n)
\]
% archive sheet number: 249
